% !TeX root = ../Fundamental-Physics-B.tex
% 第 2.1 节：牛顿以前的力学

\chapter{牛顿力学的基本定律}

\section{牛顿以前的力学}

动力学研究的是“物体为什么改变运动状态”。与第一章只描述位置、速度和加速度不同，本章要把运动状态的改变与物体所受的力联系起来。牛顿力学的基本思想可以概括为：\textbf{力不是维持运动的原因，而是改变运动状态、产生加速度的原因。}

\subsection{从日常经验到理想化模型}

从日常经验看，推箱子时手一停箱子就会停下，拉车时停止用力车也会逐渐停下，因此容易产生“力维持运动”的直觉。但这些现象中同时存在摩擦力和空气阻力。实际需要外力长期作用的，并不是“运动本身”，而是在阻力存在时仍保持某种运动状态。

伽利略通过理想斜面实验揭示了惯性的思想：若摩擦可以忽略，小球沿一个斜面滚下后，会在另一侧斜面上升到接近原来的高度；当另一侧斜面的倾角逐渐减小时，小球为了达到原来的高度会运动得越来越远。极限情况下，若另一侧变成光滑水平面，小球将一直运动下去。

\begin{knowledge}{力学观念的转变}{ch2:pre-newton-core}
\[
    \text{旧观念：力维持运动}
    \qquad \Longrightarrow \qquad
    \text{牛顿力学：力改变运动状态}
\]
物体不受外力或合外力为零时，速度不需要外力来维持；只有速度大小或方向发生变化时，才需要由合外力提供加速度。
\end{knowledge}

\subsection{惯性思想与牛顿力学的出发点}

惯性是物体保持原有运动状态的性质。一个物体若原来静止，在没有外力改变其运动状态时将继续静止；若原来做匀速直线运动，在没有外力改变其运动状态时将继续保持该匀速直线运动。

这里必须注意，静止和匀速直线运动并不是两类本质不同的状态。它们都对应于
\[
    \boldsymbol a=\boldsymbol 0.
\]
不同参考系下，同一个物体可能被描述为静止，也可能被描述为匀速直线运动。因此，讨论力学规律前必须先说明参考系。牛顿力学中，牛顿第一定律成立的参考系称为惯性参考系。

\subsection{开普勒定律的历史作用}

在牛顿力学形成以前，开普勒行星运动定律给出了行星运动的精确经验规律。这些规律后来成为牛顿建立万有引力定律的重要基础。本节不展开行星轨道推导，只需认识到：牛顿力学的建立，并不是单个公式的出现，而是从“描述运动”走向“解释运动改变原因”的理论转变。
